% arara: xelatex
\documentclass[12pt]{article}

\usepackage{physics}
\usepackage{microtype}
\usepackage{array}
\usepackage{indentfirst}
\usepackage{sectsty}
\allsectionsfont{\centering}
\usepackage{amsmath, amsfonts, amssymb}
\usepackage{comment}
\usepackage[top=2cm, left=1.2cm, right=1.2cm, bottom=2cm]{geometry}
\usepackage{lastpage}
\usepackage{enumitem}
\usepackage{caption}
\usepackage{url}
\usepackage{fancyhdr}
\usepackage{booktabs}
\usepackage{fontspec}
\usepackage{polyglossia}

\setmainlanguage{russian}
\setotherlanguages{english}
\setmainfont{Arial}
\setlength{\headheight}{15pt}

\pagestyle{fancy}
\lhead{Временные ряды}
\chead{}
\rhead{Экзаменационные билеты}
\lfoot{}
\cfoot{DON'T PANIC}
\rfoot{\thepage/\pageref{LastPage}}
\renewcommand{\headrulewidth}{0.4pt}
\renewcommand{\footrulewidth}{0.4pt}

\AddEnumerateCounter{\asbuk}{\russian@alph}{щ}
\setlist[enumerate, 2]{label=\asbuk*),ref=\asbuk*}

\DeclareMathOperator{\Cov}{\mathbb{C}ov}
\DeclareMathOperator{\Corr}{\mathbb{C}orr}
\DeclareMathOperator{\Var}{\mathbb{V}ar}
\let\P\relax
\DeclareMathOperator{\P}{\mathbb{P}}
\DeclareMathOperator{\E}{\mathbb{E}}
\DeclareMathOperator{\plim}{plim}

\newcommand \hb{\hat{\beta}}
\newcommand \hs{\hat{\sigma}}
\newcommand \htheta{\hat{\theta}}
\newcommand \s{\sigma}
\newcommand \hy{\hat{y}}
\newcommand \hY{\hat{Y}}
\newcommand \e{\varepsilon}
\newcommand \he{\hat{\e}}
\newcommand \z{z}
\newcommand \hVar{\widehat{\Var}}
\newcommand \hCorr{\widehat{\Corr}}
\newcommand \hCov{\widehat{\Cov}}
\newcommand \cN{\mathcal{N}}
\newcommand \RR{\mathbb{R}}
\newcommand \NN{\mathbb{N}}
\newcommand{\cF}{\mathcal{F}}

\begin{document}

\section*{Экзаменационные билеты}

\begin{enumerate}

    \item \textbf{Билет 1.}
    \begin{enumerate}
        \item Аддитивная и мультипликативная ETS. Связь аддитивной модели и мультипликативной через логарифмирование. Правдоподобие ETS.
    \end{enumerate}

    \item \textbf{Билет 2.}
    \begin{enumerate}
        \item Дивергенция Кульбака-Лейблера. Критерий Акаике. Как AIC помогает приближать истинную модель? Критерии применимости AIC (при каких условиях применение AIC может привести к неверным выводам?).
    \end{enumerate}

    \item \textbf{Билет 3.}
    \begin{enumerate}
        \item Белый шум. Стационарность. Автокорреляция и частная автокорреляция. Метод Юла-Волкера.
    \end{enumerate}

    \item \textbf{Билет 4.}
    \begin{enumerate}
        \item Разница между процессом и уравнением. AR(p)-процесс. Теорема о корнях лагового полинома. Стационарность и обратимость.
    \end{enumerate}

    \item \textbf{Билет 5.}
    \begin{enumerate}
        \item Нестационарность. Порядок интеграции. Теорема Вольда. Детерминированный и стохастический тренд.
    \end{enumerate}

    \item \textbf{Билет 6.}
    \begin{enumerate}
        \item ARCH-GARCH-процесс. Стационарность. Условная и безусловная дисперсия. Прогнозирование.
    \end{enumerate}

    \item \textbf{Билет 7.}
    \begin{enumerate}
        \item Копулы. Архимедовы копулы. Эллиптические копулы. Применения к расчёту риск-метрик.
    \end{enumerate}

    \item \textbf{Билет 8.}
    \begin{enumerate}
        \item Сведение к табличной задаче. Инжиниринг переменных. Борьба с утечками. Стратегии прогнозирования.
    \end{enumerate}

    \item \textbf{Билет 9.}
    \begin{enumerate}
        \item STL, MSTL. Общий алгоритм, особенности настройки. Применение к прогнозированию.
    \end{enumerate}

    \item \textbf{Билет 10.}
    \begin{enumerate}
        \item RNN, LSTM, GRU. Эволюция архитектур, особенности, преимущества и недостатки.
    \end{enumerate}

\end{enumerate}

\end{document}
